\documentclass[12pt]{article}
\usepackage[utf8]{inputenc}
\usepackage[T1]{fontenc}
\usepackage[margin=1in]{geometry}
\usepackage{setspace}
\usepackage{xcolor}
\usepackage[colorlinks=true, urlcolor=blue!60!black, linkcolor=blue!60!black]{hyperref}
\setlength{\parindent}{2em}
\setlength{\parskip}{1.25em}
\renewcommand{\baselinestretch}{1.0}
\begin{document}
\begin{center}
{\LARGE \textit{Trump v. Slaughter}, 609 U.S. 422 (2026)\par}
\vspace{0.8em}
{\large Case Brief --- Four-Page Edition\par}
\end{center}
\vspace{1em}

\section{Memory Jogger}

By 6--3, the Court overruled \textit{Humphrey's Executor v. United States} (1935) and held that the FTC's ``for-cause'' removal protections violate Article II, because principal officers who exercise executive power must be removable by the President at will (\href{https://supreme.justia.com/cases/federal/us/609/25-332/}{Justia}).

\section{Detailed Case Facts}

The Federal Trade Commission enforces more than 80 federal statutes. It issues substantive rules with the force of law, investigates businesses, conducts in-house enforcement adjudications, and sues in federal court for injunctions, redress, and civil penalties. Its five Commissioners serve staggered seven-year terms, no more than three from one party, and by statute may be removed only ``for inefficiency, neglect of duty, or malfeasance in office.''

Democrat Rebecca Kelly Slaughter was confirmed in 2018 on President Trump's nomination and reconfirmed in 2024. In March 2025, President Trump fired her and fellow Democrat Alvaro Bedoya without asserting any statutory cause; the letters said their ``continued service on the FTC [was] inconsistent with [his] Administration's priorities'' and invoked his ``authority under Article II of the Constitution.'' Bedoya later resigned, leaving Slaughter's claims alone.

\section{Procedural History}

The district court granted Slaughter summary judgment, holding \textit{Humphrey's Executor} binding, and enjoined interference with her duties. A divided D.C. Circuit panel denied a stay, Judge Rao dissenting. Over a dissent by Justice Kagan, the Supreme Court stayed the injunction in September 2025 and granted certiorari before judgment on whether the protections are valid and whether courts may prevent removal from office. After argument on December 8, 2025, the Court reversed on June 29, 2026.

\section{Judicial Votes}

Reversed and remanded, \textbf{6--3}.
\begin{itemize}
\item \textbf{Majority:} Chief Justice Roberts, joined by Justices Alito, Gorsuch, Kavanaugh, and Barrett, and by Justice Thomas except Part III--B.
\item \textbf{Concurrence:} Justice Gorsuch.
\item \textbf{Dissent:} Justice Sotomayor, joined by Justices Kagan and Jackson.
\end{itemize}

\section{Holding}

The Vesting and Take Care Clauses give the President unrestricted power to remove principal officers who exercise executive power on his behalf. The FTC's for-cause protections violate the separation of powers, and \textit{Humphrey's Executor} is overruled to the extent anything remained of it beyond entities exercising ``no part of the executive power.'' The Court did not reach the remedial question, addressed in \textit{Trump v. Cook}. Part III--B reserved advisory and chartered bodies, the Federal Reserve, and non-Article III judges.

\section{Analysis of Opinions}

\textbf{Majority (Roberts).} The Framers chose a single President for ``vigor and activity'' and accountability, and subordinate officers act as his ``assistants or deputies,'' which requires at-will removal---a principle liquidated by the Decision of 1789, early practice, and \textit{Myers v. United States}. \textit{Humphrey's} rested on an ``almost fictional view'' that the FTC exercised ``no part of the executive power,'' and \textit{Buckley}, \textit{Bowsher}, \textit{Morrison}, \textit{Free Enterprise Fund}, and \textit{Seila Law} eroded its premises. Under the \textit{Knick} factors it was overruled: its reasoning was flawed, it conflicted with later precedent, its ``quasi'' test proved unworkable, and ``[n]o branch may rely on adverse possession to claim power that the Constitution vests elsewhere.'' The modern FTC's rulemaking, adjudication, and litigation are ``executive through and through.'' A ``reasonableness'' test would let Congress insulate even the DOJ or EPA.

\textbf{Gorsuch (concurring).} Congress delegated sweeping power to independent agencies assuming insulation from the White House; now ``the fourth branch's powers still exist; they have just been reassigned to the President.'' Given the veto's ``ratchet effect,'' the Court should ``finish the journey'' by enforcing nondelegation, major questions, vagueness, Article III, and the Seventh Amendment.

\textbf{Sotomayor (dissenting).} The majority discarded 90 years of precedent and 140 years of bipartisan practice since the 1887 Interstate Commerce Commission. The Constitution is silent on removal, the Decision of 1789 barred only Senate participation, and founding-era bodies such as the Sinking Fund Commission show tenure limits were accepted. Overruling worked a ``profound bait and switch'' on those who built dozens of bipartisan commissions.

\section{Examples: Future Applications}

\begin{enumerate}
\item \textbf{Same-Side (Removable): NLRB or FERC Member.} Both issue rules, adjudicate, and sue, placing them in the ``heartland of executive power''; at-will removal is valid.
\item \textbf{Same-Side (Removable): CPSC or NRC Commissioner.} Commissions issuing binding safety rules and enforcing mandates against private parties lose their for-cause protection.
\item \textbf{Opposite-Side (Protected): Advisory Commission.} A body that only researches and reports to Congress wields ``no part of the executive power,'' so tenure protection survives under Part III--B and \textit{Buckley}.
\item \textbf{Opposite-Side (Protected): Federal Reserve Governor.} \textit{Trump v. Cook} held the same day that the Federal Reserve is an exception, given its lineage from the First and Second Banks.
\item \textbf{Fence-Sitter: Tax Court Judges and Quorum Rules.} Removing a Tax Court judge, or firing all opposition commissioners and attacking quorum and party-balance rules, lies on the frontier Part III--B left open.
\end{enumerate}

\section{Critique}

\textbf{Unitary-executive views.} Supporters hail the end of the ``headless `fourth branch''' and its fictional ``quasi'' labels. But Aaron Nielson warns of ``severability stacking'': by severing congressional checks while leaving vast delegations intact, the Court gave the President a ``statutory windfall,'' and Andy Smarick called Gorsuch's concurrence a ``con-jority opinion.''

\textbf{Progressive and practical critiques.} Critics call the majority's history selective, fault its treatment of \textit{stare decisis}, and see inconsistency with \textit{Cook}'s Federal Reserve carveout. Former FERC Chair Joseph Kelliher warns of a ``Night of the Long Knives'' at every transition, politicizing expert commissions.

\section{Key Quotations}

\begin{quote}
``When an agency `executes' a congressional mandate against private parties, it exercises executive power---no ifs, ands, or quasis about it.''
\end{quote}
\noindent --- Chief Justice Roberts, for the Court.

\begin{quote}
``If anything more is left of \textit{Humphrey's}, we overrule it. \textit{Humphrey's} has for decades been a result in search of a rationale.''
\end{quote}
\noindent --- Chief Justice Roberts, for the Court.

\begin{quote}
``Open-ended delegations of legislative power have not gone away; now they will just be exercised by agency officials who answer to the President. The power to write new regulatory crimes still exists, but now the pen ultimately rests in the President's hand. The ability to judge disputes in-house remains, but now the house is white.''
\end{quote}
\noindent --- Justice Gorsuch, concurring.

\end{document}
